\clearpage\section{把知识用到操作、波形和提问中}
\subsection{从源码到灯亮，中间发生了什么}
\begin{enumerate}
\item 写设计文件，规定电路的输入输出、寄存器和更新规则。顶层模块负责把各部分接起来。
\item 写testbench，在电脑里提供输入，做行为仿真。它适合检查加减选择、时钟沿写入、计满更新、抖动过滤等规则。
\item 配置约束。\code{pin.xdc}把顶层端口对应到芯片封装引脚，并指定接口电平标准；时序设计还要给工具说明时钟周期。比如100MHz对应10ns。
\item 检查代码并执行综合，把RTL转换成逻辑资源和寄存器的连接。RTL图帮助理解结构，综合后电路图展示优化后的结构，二者不一定长得一样。
\item 执行实现，把资源放置到芯片具体位置并完成布线，然后检查时序是否满足目标时钟要求。
\item 生成bitstream配置文件，通过Hardware Manager下载到开发板，按照操作表检查真实现象。
\end{enumerate}

仿真正确说明测试覆盖到的逻辑行为符合预期。真实板上还涉及引脚、电平、时钟约束、机械输入与实现时序，因此仿真通过后仍需综合、实现和上板。

\subsection{XDC为什么影响实物，却不决定加减规则}
实验1里的\code{input_a[0]}是逻辑端口，约束把它接到EGO1大拨码SW0对应的R1引脚。若逻辑正确、引脚却接错，板上操作就对不上。反过来，换一个引脚并不会把加法规则变成减法，规则仍来自RTL。

仿真波形中使用的是逻辑端口名，不会因为没有连到实际拨码就不能给它赋值。仿真测试和真实板上输入是两种不同的输入来源。

\pagepart{读波形与排查问题时，先缩小到一条因果链}
\subsubsection{先确认时间和进制}
看清时间单位是ns还是ms。仿真用了缩短参数时，先用实际覆盖值计算周期。二进制\code{00010000}与十六进制\code{10}是同一组位，若当成十进制10去理解，灯的位置就会判断错误。

\subsubsection{组合逻辑先看输入，时序逻辑先找时钟沿}
实验1从输入变化追到输出。实验2先找到写入沿，再查看沿前的使能和数据。实验3先找到cnt到上限的一拍，再看下一更新沿的cnt和led。实验4则把路径缩小到\code{button -> clean -> press -> count}，逐级核对。

多级寄存器不要期待所有信号在同一时刻改变。前级刚产生的值通常要到下一时钟沿才会被下游时序逻辑接收。查看波形时给变化点适当放大，分清沿前值、沿后值以及组合传播。

\begin{center}
\begin{tabularx}{\linewidth}{@{}XX@{}}\toprule
看到的现象 & 优先检查什么\\\midrule
实验1四个灯一直亮 & 是否仍在en=0；是否误拨了大拨码而非使能小拨码\\
实验2q不随d改变 & 当前是否有上升沿、en是否为1；寄存器堆是否读取写入地址\\
实验3一直停在LED0 & 是否按住复位；时钟是否连对；是否误用巨大仿真上限\\
实验3快档看起来多个灯亮 & 先回慢档观察；快档的视觉现象不直接等于逻辑图案\\
实验4单个数字不对或缺段 & 段码位序、有效电平、左右段选与引脚约束\\
实验4数字对，但位置颠倒 & 位选与display切片的对应关系\\
实验4一按增加很多次 & 当前看的是RR还是DD；消抖参数是否仍是仿真值\\
实验4暂停后数字仍闪或熄灭 & 是否误停了扫描；当前设计只暂停TT计时\\
波形出现X & 先确认复位、端口连接与是否有多个驱动\\\bottomrule
\end{tabularx}
\end{center}

\subsection{代码没有写CPU式循环，扫描为何一直重复}
时钟不断提供上升沿。周期计数器在每拍更新，计满时扫描编号加一，三位编号溢出后回0。各模块从配置后就持续工作，因此无需写软件式\code{while}循环。硬件计数器的状态转移本身就形成重复行为。

\pagepart{自测题：先自己推演，再看下一页}
\begin{enumerate}[label=\arabic*.]
\item 四位无符号数最大是多少？\code{8'h81}接八个LED，哪些灯亮？
\item 实验1中en=0、a=5、b=2，改变选择开关会改变输出吗？如果en=1、a=15、b=1并选加法呢？
\item D触发器clr=0、en=1，d在两个上升沿之间由0变1，q什么时候更新？en=0时呢？
\item 寄存器堆的2号地址已有A5，5号地址已有3C。en=0时，把rsel从2改成5，q会变吗？为什么？
\item 100MHz下每0.1秒更新一次，为什么上限是9,999,999？若上限设为6，每次更新相隔几拍？
\item 实验3慢档运行到cnt=5,000,000，直接切到快档，\code{>=}为什么比\code{==}合适？
\item \code{led_pos=8'h80}时，按编号递增方向的下一个位置是什么？若下一位置是80，对应双灯图案是什么？
\item 边沿检测器原来r0=r1=0，输入已稳定为1。第一拍之后r0、r1各是多少？第二拍之后呢？
\item 消抖后按键高电平持续一秒，为何仍需要边沿检测？两级同步能直接替代20ms消抖吗？
\item 每位显示1ms，八位一轮是多少ms、整轮刷新率多少Hz？扫描编号0选中最左还是最右？
\item \code{32'h25310489}在这里保存什么格式的数据？把普通整数20的二进制\code{00010100}直接拆成两个四位数字，会显示什么？
\item TT从00至20每0.1秒变化一次，一轮多久？按S2暂停时，数码管扫描是否也停止？
\end{enumerate}

\subsection{试着完整讲一次S3按下的过程}
从手指按下开始，沿着“机械触点变化、输入采样、稳定判断、边沿通知、BCD加一、显示数据、扫描、段码”逐项解释。RR支路绕过哪一步？哪几步需要等待时钟？屏幕上看到新数字还会受当前扫描位置怎样影响？


\pagepart{自测参考与一页复习索引}
\begin{enumerate}[label=\arabic*.]
\item 最大15。81是\code{10000001}，第7、0位灯亮。
\item 禁用时始终1111，选择开关不影响结果。15+1的低四位为0000。
\item 下一次时钟上升沿且en=1时写入；en=0时保持原值。clr有效会优先异步清零。
\item 会从A5变为3C，因为读出是组合选择。en只控制写入。
\item 0到9,999,999有1000万个计数状态，每拍10ns，共0.1s。0到6有7个状态，更新间隔7拍。
\item 旧计数已经超过199,999，\code{>=}能及时识别超出；\code{==}可能错过新的上限。
\item 80循环左移后为01；位置80对应的双灯是81。
\item 第一拍后r0=1、r1=0，通知为1；第二拍后均为1，通知为0。下游计数器在第二拍接收到沿前有效的通知。
\item 高电平持续很多拍，电平作使能会重复计数；边沿检测仅给一次通知。两级同步的时间尺度很短，不能代替毫秒级消抖。
\item 一轮8ms，125Hz。编号0选择最右位。
\item 每四位装一个待显示的数字。\code{00010100}拆开是1与4，直接译码会显示14；BCD的20应写成\code{00100000}。
\item 共21个值，一轮2.1s。S2暂停TT计时，扫描继续；RR和DD也不受running控制。
\end{enumerate}

\begin{center}
\begin{tabularx}{\linewidth}{@{}lX@{}}\toprule
术语 & 在本实验中的具体意义\\\midrule
组合逻辑 & 当前输入决定输出，例如加减选择、段码译码\\
时序逻辑 & 寄存器保存状态，并按时钟与复位规则更新\\
寄存器 & 保存一组位，例如cnt、led\_pos、ones\\
使能 & 允许本拍更新，通常写在if条件里\\
参数 & 展开设计时确定的配置；testbench可覆盖为小数值\\
同步 & 对来自时钟域外的输入采样，降低亚稳态传播风险\\
消抖 & 稳定持续达到窗口才接受新电平\\
边沿检测 & 比较相邻采样，将上升变化变成单拍通知\\
BCD & 每个十进制数字单独占四位\\
位选／段选 & 选择显示位置／决定该位置的字形\\\bottomrule
\end{tabularx}
\end{center}
